% TOP_PROBLEM: {"id": 2638, "title": "Kourovka Notebook Problem 21.129", "category_id": 17, "category": {"id": 17, "name": "group_theory", "display_name": "Group Theory", "description": "Problems about groups, group actions, representations, and related algebraic structures.", "slug": "group-theory", "order_index": 17, "created_at": "2026-07-31T15:26:25.671Z"}, "status": "open", "problem_number": "KOU-21.129", "set_id": 10}
% FIRST_POSTED: 2026-10-01
% ENTRY_KIND: statement_only
% UnsolvedMath ID 2638; source catalogue code KOU-21.129
% !TeX program = lualatex
% Definition and sources only; this file is not an LLM solution attempt.
\documentclass[11pt,a4paper]{article}
\usepackage[margin=25mm]{geometry}
\usepackage{fontspec}
\setmainfont{lmroman10-regular.otf}[BoldFont=lmroman10-bold.otf,
 ItalicFont=lmroman10-italic.otf,BoldItalicFont=lmroman10-bolditalic.otf]
\usepackage{amsmath,amssymb,mathtools}
\usepackage{unicode-math}
\setmathfont{latinmodern-math.otf}
\AtBeginDocument{\let\setminus\smallsetminus}
\usepackage{xurl,hyperref}
\hypersetup{colorlinks=true,urlcolor=blue}
\setcounter{secnumdepth}{0}
\setlength{\parindent}{0pt}
\setlength{\parskip}{5pt}
\setlength{\emergencystretch}{3em}
\begin{document}
\section*{Kourovka Notebook Problem 21.129}\label{2638}
{\small UnsolvedMath ID 2638 · KOU-21.129 · Group Theory · Kourovka Notebook - New Problems, Issue 21}
\subsection{Definitions and mathematical statement}
A Coxeter matrix is a finite symmetric matrix $M=(m_{ij})$ with $m_{ii}=1$ and $m_{ij}\in\{2,3,\ldots\}\cup\{\infty\}$ for $i\ne j$. Let $[s_i,s_j]_m$ mean the alternating word of length $m$ starting with $s_i$. The associated Artin group is
\[
 A_M=\langle s_1,\ldots,s_r\mid
 [s_i,s_j]_{m_{ij}}=[s_j,s_i]_{m_{ij}}
 \text{ whenever }i<j,\ m_{ij}<\infty\rangle.
\]
Its Coxeter group is obtained by additionally imposing $s_i^2=1$ for all $i$. The Artin group is of spherical type when this Coxeter group is finite.

Equip each finitely generated group with a word metric. A quasi-isometry from $G$ to $H$ is a map $f:G\to H$ with constants $\lambda\ge1$, $C\ge0$ such that
\[
 \lambda^{-1}d_G(x,y)-C\le d_H(f(x),f(y))
 \le\lambda d_G(x,y)+C
\]
for all $x,y$, and every point of $H$ lies within distance $C$ of $f(G)$. This existence property does not depend on the finite generating sets chosen.

If two spherical-type Artin groups are quasi-isometric, must they be commensurable? Here commensurable means that some finite-index subgroup of one is isomorphic to some finite-index subgroup of the other, with no requirement that the indices agree. The groups may come from different matrices and need not have irreducible Coxeter diagrams. The proposed conclusion is algebraic despite the metric nature of the hypothesis.
\subsection{Short English statement}
Must quasi-isometric Artin groups of spherical type have isomorphic finite-index subgroups?
\subsection{Sources}
\begin{itemize}
\item[S1] ULAM.AI and the UnsolvedMath Contributors, supplied problems.json, record 2638 (KOU-21.129), Kourovka Notebook - New Problems, Issue 21. Catalogue identity and source transcription; CC BY 4.0.
\url{https://huggingface.co/datasets/ulamai/UnsolvedMath}.
\item[S2] The Kourovka Notebook, 21st edition, version 47 (30 September 2026), new problems, Problem 21.129. Expanded formulation of the supplied catalogue statement.
\url{https://arxiv.org/pdf/1401.0300v47}.
\end{itemize}
\end{document}
